\section{Упругие (механические) волны в твёрдых телах}

\subsection{Уравнения Ламе}

В изотропной линейно-упругой твёрдой среде динамика малых перемещений
$\mathbf{u}(x,t)$ описывается уравнением Ламе (или уравнением Навье):
\begin{equation}
\rho\,\frac{\partial^2 \mathbf{u}}{\partial t^2}
=
(\lambda + \mu)\nabla(\nabla\cdot\mathbf{u})
+
\mu\,\nabla^2\mathbf{u},
\label{eq:navier}
\end{equation}
где $\lambda$ и $\mu$ — параметры Ламе, $\rho$ — плотность.  

Это уравнение получается как комбинация закона Гука и уравнения движения
(закона сохранения импульса) механики сплошной среды \cite{MIT_continuum, Berkeley_ME185}.

\subsection{Типы упругих волн}

Из уравнения~(\ref{eq:navier}) следуют следующие волны:

\paragraph{P-волны (продольные)} — смещения параллельны направлению распространения.  
\paragraph{S-волны (поперечные)} — смещения ортогональны направлению распространения.  
\paragraph{Поверхностные волны Релея} — волны, бегущие вдоль свободной поверхности твёрдого тела, затухающие с глубиной \cite{Achenbach, Graff}.  
\paragraph{Волны Лява} — специфичный тип волн, возникающих в слоистых анизотропных средах (используется в геофизике).  

\subsection{Скорости волн}

\paragraph{Продольные:}
\begin{equation}
v_P = \sqrt{\frac{\lambda + 2\mu}{\rho}}.
\end{equation}

\paragraph{Поперечные:}
\begin{equation}
v_S = \sqrt{\frac{\mu}{\rho}}.
\end{equation}

Поскольку $\mu > 0$, следует:
\[
v_S < v_P.
\]

\subsection{Отражение и преломление (базовое геофизическое приближение)}

При падении упругой волны на границу двух упругих сред часть энергии отражается, часть преломляется.  
Коэффициенты отражения/передачи связаны с механическими импедансами
\[
Z_P = \rho\, v_P,\qquad Z_S = \rho\, v_S.
\]
При наклонном падении волны возможен переход энергии между типами волн (P → P, P → S, и т.д.), что лежит в основе геофизической интерпретации сейсмических записей \cite{Achenbach, Berkeley_ME185}.



\section*{Список источников}

\begin{thebibliography}{9}

\bibitem{MIT_continuum}
MIT OpenCourseWare. \textit{Applications of Continuum Mechanics to Earth, Atmospheric and Planetary Sciences}. Lecture notes (2006).  
URL: \url{https://ocw.mit.edu/courses/12-005-applications-of-continuum-mechanics-to-earth-atmospheric-and-planetary-sciences-spring-2006/lecture-notes/}

\bibitem{Berkeley_ME185}
University of California, Berkeley. \textit{ME185 — Introduction to Continuum Mechanics (lecture notes)}.  
URL: \url{https://csml.berkeley.edu/Notes/ME185.pdf}

\bibitem{Achenbach}
Achenbach, J. D. \textit{Wave Propagation in Elastic Solids}. Dover Publications (репринт).  
PDF-версия: \url{https://archive.org/details/wavepropagationinelasticsolids}

\bibitem{Graff}
Graff, K. F. \textit{Wave Motion in Elastic Solids}. Dover Publications.  
PDF-версия: \url{https://archive.org/details/wavemotioninelasticsolids}

\end{thebibliography}
